\documentclass[journal]{IEEEtran}


%% 1. Geometry and Margins
\usepackage[a4paper,top=3cm,bottom=2cm,left=3cm,right=3cm,marginparwidth=2cm]{geometry}

%%%%%%%%%%%%%%%%%%%%%%%%%%%%%% Package Imports %%%%%%%%%%%%%%%%%%%%%%%%%%%%%%
\usepackage[english]{babel}
\usepackage[T1]{fontenc}
% Removed inputenc (not needed in 2024+)
\usepackage{graphicx}
\usepackage[normalem]{ulem}
\usepackage[colorinlistoftodos]{todonotes}
\usepackage{setspace}
\usepackage{amsmath, amssymb} % Added amssymb for extra symbols
\usepackage{bm}

%% 2. Tables and Figures
\usepackage{booktabs}
\usepackage{tabularx}
\usepackage{tabularray}
\usepackage{float}
\usepackage{hyperref}


%% 3. Styling and Fonts
\usepackage{xcolor}
\usepackage{listings}
\usepackage{titling}
\usepackage{authblk}
\usepackage{siunitx}
\usepackage{nicefrac}
\usepackage{blindtext}

%% 4. Colors (TU Delft palette)
\definecolor{tudelftdarkblue}{RGB}{12,35,64}
\definecolor{tudelftcyan}{RGB}{0,166,214}
\definecolor{tudelftblue}{RGB}{0, 118, 194}
\definecolor{titleblue}{rgb}{0.2, 0.4, 0.6}
\definecolor{backcolour}{rgb}{0.92,0.94,0.96}
\hyphenation{op-tical net-works semi-conduc-tor}

% Konfiguriere hyperref, um die Schrift zu färben statt Kästen zu zeichnen
\hypersetup{
    colorlinks=true,    % Schaltet die Färbung der Schrift ein
    linkcolor=tudelftblue, % Farbe für interne Links (Figures, Sections, etc.)
    citecolor=tudelftblue, % Farbe für Zitate (References)
    urlcolor=tudelftblue   % Farbe für Weblinks
}

\begin{document}
%
% paper title
% Titles are generally capitalized except for words such as a, an, and, as,
% at, but, by, for, in, nor, of, on, or, the, to and up, which are usually
% not capitalized unless they are the first or last word of the title.
% Linebreaks \\ can be used within to get better formatting as desired.
% Do not put math or special symbols in the title.
\title{How Do Modern High-Voltage DC-DC Converter Topologies (e.g., Resonant vs. Hardswitched) Compare in Terms of Efficiency and Mass for Ion Thruster PPUs?}


\author{Paul Schneider,~\IEEEmembership{Student,~EPE,}
        Giorgio Recchilongo,~\IEEEmembership{Student,~ME,}
        and~Felix~Voßbrink,~\IEEEmembership{Student,~EPE}% <-this % stops a space
\thanks{This article was written as part of the course EE4C01 Profile Orientation \& Academic Skills of the EEMCS Faculty at TU Delft. Total word count: 2099}}

% note the % following the last \IEEEmembership and also \thanks - 
% these prevent an unwanted space from occurring between the last author name
% and the end of the author line. i.e., if you had this:
% 
% \author{....lastname \thanks{...} \thanks{...} }
%                     ^------------^------------^----Do not want these spaces!
%
% a space would be appended to the last name and could cause every name on that
% line to be shifted left slightly. This is one of those "LaTeX things". For
% instance, "\textbf{A} \textbf{B}" will typeset as "A B" not "AB". To get
% "AB" then you have to do: "\textbf{A}\textbf{B}"
% \thanks is no different in this regard, so shield the last } of each \thanks
% that ends a line with a % and do not let a space in before the next \thanks.
% Spaces after \IEEEmembership other than the last one are OK (and needed) as
% you are supposed to have spaces between the names. For what it is worth,
% this is a minor point as most people would not even notice if the said evil
% space somehow managed to creep in.



% The paper headers
% \markboth{}%
% {Shell \MakeLowercase{\textit{et al.}}: Bare Demo of IEEEtran.cls for IEEE Journals}










% make the title area
\maketitle

% As a general rule, do not put math, special symbols or citations
% in the abstract or keywords.
\begin{abstract}
Ion thrusters have emerged as a primary propulsion choice for long-duration space missions. Their performance is heavily dependent on the Power Processing Unit (PPU) and its ability to efficiently step up low bus voltages to kilovolt levels. This paper evaluates the trade-offs between hard-switched and soft-switched (resonant) DC-DC converter topologies in terms of efficiency, mass, and reliability in the context of emerging power electronic semiconductors, such as gallium arsenide (GaAs) and silicon carbide (SiC).

While hard-switched architectures inherently suffer from their switching losses being proportional to the operating frequency, and thus limited power density, resonant topologies utilize Zero Voltage Switching (ZVS) and Zero Current Switching (ZCS) to decouple efficiency from frequency. Our review finds that while soft-switching can achieve efficiencies up to 99~\% and reduce electromagnetic interference, it introduces additional mass by advanced circuit complexity. The analysis concludes that modern PPU design must balance the superior high-frequency efficiency of resonant converters against the increased component count and potential reliability vulnerabilities inherent in the complex architectures.
\end{abstract}

% Note that keywords are not normally used for peerreview papers.
% \begin{IEEEkeywords}
% IEEE, IEEEtran, journal, \LaTeX, paper, template.
% \end{IEEEkeywords}


\IEEEpeerreviewmaketitle


\section{Introduction}
%% Giorgio

%% Paul
% Hard Switched DC-DC Converters: What they are, why they were used, and what are their problems regarding switching losses and EMI.
% Softswitched DC-DC Converters: What they are, how they work, and their advantages in terms of reduced switching losses and EMI. But they come with increased complexity and therefore weight, which is crucial in satellite applications.
% Something like this (gemini): 
For the DC DC converter different topologies exist, which can be broadly categorized into hard-switched and soft-switched (resonant) converters. 
The primary motivation for investigating resonant converter topologies in ion propulsion is the relationship between switching frequency and converter system mass. In traditional power electronic circuits, using hard-switched devices, switching losses are directly proportional to the operating frequency. While increasing frequency is the standard method for reducing the size of magnetic components like inductors, and therefore their weight, it leads to power dissipation in hard-switched systems where voltage and current overlap during transitions.Theoretical research~\cite{bellar1998review}, suggests that by integrating a resonant network, the converter can achieve Zero Voltage Switching (ZVS) or Zero Current Switching (ZCS). This effectively shapes the switching waveforms so that transitions occur only when the voltage or current is zero, hence the power disspated during switching is minimized. The key advantage is hence, decoupling efficiency from frequency. This is especially relevant for converters for satellite applications, such as ion thrusters, where the core engineering challenge is maximizing the effective specific power ($kW/kg$) of components~\cite{gao2017soft}. On the other hand, soft-switched converters need additional circuitry to achieve ZVS or ZCS, which increases the overall system complexity and mass. Hence a comparison of the two topologies in terms of efficiency and mass is crucial for determining the optimal design for ion thruster PPUs.
\section{Methods}
% First step. Ai deep research for overview. From there keyword searches filtering for specific thruster architectures to compare their DCDC converter topologies
To conduct a comprehensive comparison of modern high-voltage DC-DC converter topologies for ion thruster Power Processing Units (PPUs), we employed a systematic research methodology. We started by performing an extensive literature review to gain a broad understanding of the various DC-DC converter topologies, utilising tools such as deep research through LLM-based search engines. From this, we went on to conduct keyword searches in scientific databases, filtering for specific thruster architectures and their associated DC-DC converter topologies. This allowed us to identify relevant studies and technical papers that provided insights into the efficiency and mass characteristics of resonant and hard-switched converter topologies.


\section{Topology Comparison}
\subsection{Hard Switching}
\subsection{Soft Switching: Resonant Converters}
\subsubsection{Zero Voltage Switching}
\subsubsection{Zero Current Switching}

\subsection{Efficiency Comparison}
\subsection{Weight Comparison}
The mass of ion thruster power conversion units is a critical constraint in space applications, where payload weight directly dictates launch costs and mission feasibility. Hard-switching architectures necessitate heavy external dv/dt filters to suppress voltage spikes and protect sensitive components. As noted by Cai et al. \cite{cai2019voltage}, these passive filters significantly increase the overall size and weight of the propulsion system. While soft-switching can reduce total mass by integrating the filtering effect into the switching transition, research into weight optimization through the control of these transitions remains limited.

\subsection{Reliability Comparison}
The reliability of power conversion for ion thrusters involves a trade-off between electrical stress and architectural complexity. Tang et al. \cite{tang1998emi} demonstrate that soft-switching improves reliability by reducing $dv/dt$ and $di/dt$ rates, which lowers electromagnetic interference (EMI) and common-mode noise compared to hard-switching. In the context of space propulsion, this reduction is critical as high EMI can interfere with sensitive spacecraft antennas and communication links. By using snubber capacitors to slow voltage transitions, soft-switching mitigates stress on insulation and parasitic components within the thruster's power processing unit. However, this is offset by a higher component count, including auxiliary switches and resonant inductors, which statistically increases potential failure points. % include paper on hard and soft switch to show the diff in component count
Furthermore, the increased system complexity required for precise timing control adds vulnerability; any control failure can lead to hard-switching under high-stress conditions, potentially negating the reliability benefits for long-duration missions.

\section{Conclusion}
The goal of this literature review was to evaluate how modern high-voltage DC-DC converter topologies, specifically resonant and hard-switched architectures, compare in terms of efficiency, mass, complexity and reliability for ion thruster PPUs. Typically, the selection between these topologies involves a trade-off between power consumption and overall system complexity. The main limitation of hard-switched architectures, normally employing full-bridge inverters and PWM blocks, is constituted by the proportionality between the operating frequency and the switching losses. In fact, even though an increase in frequency leads to a reduction of the mass of magnetic components, it also affects negatively power dissipation due to voltage-current overlap during transitions.

On the other hand, soft-switching topologies solve this issue by integrating resonant networks that implement either Zero Voltage Switching or Zero Current Switching. This effectively modifies the switching waveforms so that transitions occur only when the voltage or current is zero, thus minimizing power dissipation and making it independent of operating frequency. Such decoupling is advantageous for satellite and space applications where maximizing the metric of specific power ($kW/kg$) constitutes the main challenge. However, these benefits are achieved only by the use of additional circuitry, which increases system complexity and mass.

Overall, design for modern PPUs has to strike a delicate balance between the superior efficiency at higher frequencies of resonant converters and their additional mass. Furthermore, current trends focusing on wide-bandgap switches, using materials such as Gallium Nitride and Silicon Carbide, and the continuous miniaturization of components are leading to improvement in power density metrics. Future research should focus more on the integration of those wide-bandgap materials and newly available technologies in converters, in addition to the analysis of long-term reliability and radiation hardening of these switching components when used in a space environment. These proposed studies could improve and promote the feasibility of human deep-space exploration.

% Bibliography (Uncomment when needed)
\addcontentsline{toc}{section}{References}
\bibliographystyle{IEEEtran}
\bibliography{97_Reference/references.bib}

\end{document}


